% problem-000239 10 硫化学与亚砜手性
\section*{10. 硫化学与亚砜手性}
\textit{下列为分问。}

\subsection*{10-1}
如下所示的埃索美拉唑是一种胃肠类药物，用 $R / S$ 标注其绝对构型：

（图：）

\subsection*{10-2}
手性亚砜具有相对稳定性，但烯丙基取代的亚砜易发生外消旋化：

（图：）

此外消旋化过程与取代基和溶剂极性紧密相关。请回答以下问题：

10-2-1 以下哪个手性亚砜更易发生外消旋化？（）
(a) Ar 为 4-三氟甲基苯基；(b) 4-甲氧基苯基；(c) Ar 为苯基；(d) 无法判断。

10-2-2 为以上选择提供合理的解释: ( )
(a) 4-三氟甲基苯基为吸电子基团, 使亚砜不稳定; (b) 4-三氟甲基苯基为吸电子基团, 使亚砜稳定; (c) 4-甲氧基苯基为给电子基团, 使亚砜不稳定; (d) 4-甲氧基苯基为给电子基团, 使亚砜稳定; (e) 给电子基团和吸电子基团都使亚砜不稳定; (f) 给电子基团和吸电子基团都使亚砜稳定; (g) 无法判断。

10-2-3 溶剂极性增加，对亚砜外消旋化是有利的还是不利的？( )
(a) 有利； (b) 不利。

\subsection*{10-2}
-4 为以上选择提供合理的解释：（ ）。

(a) 次磺酸酯的极性比亚砜强, 极性溶剂更容易溶剂化亚砜; (b) 次磺酸酯的极性比亚砜强, 极性溶剂更容易溶剂化次磺酸酯; (c) 亚砜的极性比次磺酸酯强, 极性溶剂更容易溶剂化亚砜; (d) 亚砜的极性比次磺酸酯强, 极性溶剂更容易溶剂化次磺酸酯。

\subsection*{10-3}
以(S)-丁-3-烯-2-醇为原料，不对称中心可以从碳中心转移到硫中心，转化过程如下：

（图：）

\subsection*{10-3}
-1 确定 C 中碳碳双键的立体构型。

\subsection*{10-3}
-2 画出 B 的立体结构式。

\subsection*{10-3}
-3 画出由 B 转为 C 的过渡态结构。

\subsection*{10-4}
如下所示的双环[2.2.1]化合物 E 在二氯甲烷中回流转化为化合物 F。F 在乙醇回流转化为双环[2.2.1]化合物 G (需要关注立体化学)。

（图：）

\subsection*{10-4}
-1 画出 F 的立体结构式，并在其结构中用 R/S 标出每个手性中心的绝对构型。

\subsection*{10-4}
-2 画出由 F 到 G 的关键中间体。

\subsection*{10-4}
-3 判断 G 中三个手性碳原子的绝对构型，并解释由 F 到 G 转换过程中的区域选择性。

本试题和答案版权属中国化学会所有，未经中国化学会化学竞赛负责人授权，任何人不得翻印，不得在出版物或互联网网站上转载、贩卖、赢利，违者必究。本试题和相应答案将分别于2024年9月1日14:00和9月8日14:00在www.chemsoc.org.cn网站上公布
